\chapter{Оптическая решётка}\label{app:A}

%Для атома, взаимодействующего с переменным электрическим полем  $\mathrm{E} (t) = \mathrm{E}_0 \cos(\omega t)$ задаётся формулой \ref{eq:StarkShift}, в которой
%
%\begin{equation}
%	\label{eq:StarkShift}
%	\Delta E_i (\omega) = \frac{1}{2h}\sum\frac{\omega_{if} |d_{if}|^2}{\omega_{if} - \omega^2}\mathrm{E} 
%\end{equation}

Здесь будет сборник фактов и формул об оптических решётках, которые могут мне понадобиться в других местах.

\section{Оптическая решётка и режим Лэмба-Дике}

Оптическая решётка - потенциал, формируемый стоячей волной лазерного излучения.
Для вектора k сонаправленного оси Z он записывается следующим образом.
\begin{equation}
	\label{eq:LatticePotential1_app} 
	U(r, z) = -\alpha I e^{-2r^2/\omega^2(z)} \cos^2(kz)
\end{equation}

Пространственное распределение эквипотенциальных поверхностей оптической решетки приведено на рисунке. 
Движение атомов в потенциале оптической решетки в продольном направлении ограничено её периодом, то есть атомы оказываются локализованными в пространстве на размерах порядка половины длины волны формирующего оптическую решетку излучения. 
Локализация атомов в поперечном направлении определяется радиусом гауссового пучка.

Если записать гамильтониан этого безобразия, то можно получить, что существует параметр Лэмба-Дике, который подвляет переходы с изменением колебательного числа и эффект отдачи.
Получаем его из Гаммильтониана.
Задаётся такой вот формулой: $\eta = k\sqrt{\frac{\hbar}{2m\Omega_L}}$, в которой понятно что есть что.

Сдвиг энергий уровней, связанный с эффектом Штарка приводит к тому, что смещается и частота часового перехода.
При этом, зависимость сдвига энергии разных уровней от частоты излучения решётки разная, что позволяет подобрать длину волны решётки так, чтобы сдвиг, связанный с электродипольной поляризуемостью оказался бы равен нулю. 
Такая длина волны назвается магической.

В работе \cite{} был проведён поиск магических длин волн для часового перехода в атомах тулия. 
Оказалось, что длина волны $1063.5\unit{нм}$ близка к магической, и последнее время мы работаем вблизи неё. 

Подход, описанный выше, перестаёт работать при приближении к уровню неточности $10^{-17}$. 
В этот момент нужно учитывать влияние поляризуемостей старших порядков, как то магнито-дипольная электро-квадрупольная и гиперполяризуемость.
Потенциал с учётом этого всего записывается так \cite{ushijima2018operational}:  

\begin{equation}
	\label{eq:LatticePotential2_app} 
	U(r, z) = -\alpha E^2 (r) \cos^2(kz) - B^2 \sin^2 (kz) - E^2 (r) \sin^2 (kz) _ ...  
\end{equation}

А сдвиг энергии при этом, зависит не только от параметров (значений поляризуемости), но и от движения атомов в оптической решётке.
Величина этого сдвига задаётся формулой \ref{eq:LightShift}, в которой $h$ - постоянная Планка, $k_B$ - постоянная Больцмана, $u$ - глубина потенциала оптической решётки, $\delta_L$ - отстройка излучения оптической решётки от магической длины волны, $\bar{n}_z$ и $T_r$ - среднее колебательное число опрашиваемого ансамбля атомов и их радиальная температура.
Параметрами в этой формулы являются склон разностной электродипольной поляризуемости вблизи магической длины волны $\frac{\partial \alpha^{E1}}{\partial\nu}$, разностные электро-квадрупольная и магнито-дипольная поляризуемости $\alpha^{qm}$, а также разностная гиперполяризуемость $\beta$.
Характерный вид зависимостей такого сдвига от глубины решётки для предварительно-измеренных параметров \ref{tab:LightShiftsParameters} задатся на рисунке \ref{fig:LightShifts}  

\begin{table}[]
	\centering
	\begin{tabular}{|cc|}
		\hline
		Параметр & Значение\\
		\hline
		Магическая частота (ТГц)																		&  282\\
		$\frac{\partial \alpha^{E1}}{\partial\nu}$ ($h\times\unit{Гц}/(E_{rec}\times\unit{ГГц})$)		&  1\\
		$\alpha^{qm}$ () 																				&  2\\
		$\beta$ ()																						&  3\\
		\hline
	\end{tabular}
	\caption{Параметры, отвечающие за световой сдвиг в атомах тулия. Известны предварительно.}
	\label{tab:LightShiftsParameters}
\end{table}

\begin{equation}
	\label{eq:LightShift}
	\begin{multlined}
		h\Delta\nu_{LS} (u, \delta_L, \bar{n}_Z, T_r) \approx (\frac{\partial \alpha^{E1}}{\partial\nu} \delta_L - \alpha^{qm})(\bar{n}_z + \frac{1}{2})u^{1/2} - \\
        [\frac{\partial \alpha^{E1}}{\partial\nu} \delta_L + \frac{3}{2}\beta(\bar{n}_z^2 + \bar{n}_z + \frac{1}{2})]u + \\
        2\beta(\bar{n}_z+\frac{1}{2})u^{3/2} - \beta u^2 , u^j \approx (1 - \frac{j k_B T_r}{uE_r})u^j
	\end{multlined}	
\end{equation}

\begin{figure}[ht]
	\centerfloat{
		\includegraphics[scale=0.8]{AppendixA/LightShifts}
	}
	\caption{
		Индуцированные оптической решёткой сдвиги частоты часового перехода для различных параметров эксперимента.		
	}\label{fig:LightShifts}
\end{figure} 
	